\section{Counterfactual Effect Binding}
\label{sec:method}

\begin{figure*}[t]
\centering
\begin{tikzpicture}[
  font=\scriptsize,
  actor/.style={circle, minimum size=10mm, inner sep=0pt,
                draw=TEDark!65, fill=white, line width=.65pt},
  lane/.style={draw=TEGold!75, fill=TEGoldLight, rounded corners=5pt,
               line width=.6pt},
  message/.style={draw=TEDark!55, fill=white, rounded corners=2pt,
                  align=center, minimum height=8mm, inner sep=3pt},
  flow/.style={-{Latex[length=2.1mm]}, draw=TEDark!75, line width=.75pt},
  proposal/.style={flow, draw=TEOrange},
  evidence/.style={flow, draw=TEBlue},
  allow/.style={flow, draw=TETeal, line width=.85pt},
  stop/.style={flow, draw=TERed, line width=.85pt},
  title/.style={font=\small\bfseries, anchor=west, text=TEDark},
  label/.style={font=\scriptsize, align=center, text=TEDark!85},
  arrowlabel/.style={font=\scriptsize, fill=white, inner sep=1.3pt},
]

% Panel headings and divider.
\node[title] at (-8.25,2.10) {(a) OFFLINE: REGISTER THE EFFECT INTERFACE};
\node[title] at (.35,2.10) {(b) ONLINE: MEDIATE AN EXACT CALL};
\draw[TEDark!22, line width=.5pt] (0,2.30)--(0,-2.55);

% ---------------------------------------------------------------------------
% Offline: an untrusted proposal is tested against executable behavior.
\node[actor, draw=TEBlue, fill=TEBlueLight] (toolactor) at (-7.30,.55) {};
\path (toolactor.center) pic[scale=1.0] {te-tool};
\node[label, text width=1.55cm] at (-7.30,-.18)
  {Tool schema +\\implementation};

\path[lane] (-5.60,1.40) rectangle (-2.80,-1.18);
\node[font=\scriptsize\bfseries, text=TEGold] at (-4.20,1.13)
  {Counterfactual registry};
\node[actor, minimum size=8mm, draw=TEBlue, fill=TEBlueLight]
  (sandbox) at (-4.20,.48) {};
\path (sandbox.center) pic[scale=.77] {te-flask};
\node[label] at (-4.20,.02) {execute base and mutants};
\draw[flow, draw=TEGold!80, -{Latex[length=1.7mm]}]
  (-4.20,-.14)--(-4.20,-.43);
\node[actor, minimum size=8mm, draw=TETeal, fill=TETealLight]
  (diff) at (-4.20,-.72) {};
\path (diff.center) pic[scale=.77] {te-inspect};
\node[label] at (-4.20,-1.03) {diff effects; test separation};

\node[actor, draw=TEOrange, fill=TEOrangeLight] (proposer) at (-1.12,.72) {};
\path (proposer.center) pic[scale=1.0] {te-chip};
\node[label, text width=1.58cm] at (-1.12,-.04)
  {Proposer\\developer or LLM};

\draw[evidence] (toolactor.east) -- node[arrowlabel, above, text=TEBlue]
  {executable schema} (-5.60,.55);
\draw[proposal] (proposer.west) -- node[arrowlabel, above, text=TEOrange]
  {candidate $d_t$} (-2.80,.72);
\draw[proposal] (-2.80,-.52) -- node[arrowlabel, below, text=TEOrange]
  {counterexample} (-.48,-.52) |- (proposer.east);

\node[actor, minimum size=8.5mm, draw=TETeal, fill=TETealLight]
  (frozen) at (-4.20,-1.82) {};
\path (frozen.center) pic[scale=.82] {te-lock};
\draw[allow] (-4.20,-1.18) -- (frozen.north);
\node[label, text=TETeal] at (-4.20,-2.32)
  {freeze validated $d_t$; otherwise review};

% ---------------------------------------------------------------------------
% Online: task authority and the exact proposed call meet at pre-commit.
\node[message, text width=2.10cm] (task) at (1.35,1.13)
  {\textbf{Authenticated task}\\typed authority};
\node[actor, draw=TEOrange, fill=TEOrangeLight] (agent) at (7.35,1.13) {};
\path (agent.center) pic[scale=1.0] {te-chip};
\node[label, text width=1.45cm] at (7.35,.43)
  {LLM agent\\\textit{untrusted}};
\draw[flow] (task.east) -- node[arrowlabel, above] {task} (agent.west);

\path[lane] (3.15,.58) rectangle (5.55,-1.28);
\node[font=\scriptsize\bfseries, text=TEGold] at (4.35,.31)
  {Pre-commit monitor};
\node[actor, minimum size=8mm, draw=TETeal, fill=TETealLight]
  (guardicon) at (4.35,-.28) {};
\path (guardicon.center) pic[scale=.77] {te-shield};
\node[label] at (4.35,-.82)
  {instantiate atoms\\compare with authority};

\draw[evidence] (task.south) |- node[pos=.72, arrowlabel, above, text=TEBlue]
  {authority} (3.15,.26);
\draw[proposal] (agent.south west) |- node[pos=.30, arrowlabel, below,
  text=TEOrange] {exact call $c=(t,\bar{x})$} (5.55,-.02);

\node[actor, draw=TETeal, fill=TETealLight] (toolrun) at (1.35,-1.42) {};
\path (toolrun.center) pic[scale=1.0] {te-tool};
\node[label, text=TETeal] at (1.35,-2.08) {Execute checked call};
\draw[allow] (3.15,-1.02) -- node[arrowlabel, below, text=TETeal]
  {\Allow{}} (toolrun.east);

\node[actor, minimum size=8.5mm, draw=TERed, fill=TERedLight]
  (blocked) at (7.05,-1.42) {};
\path (blocked.center) pic[scale=.82] {te-stop};
\draw[stop] (5.55,-1.02) -- node[arrowlabel, below, text=TERed]
  {\Deny{} / \Abstain} (blocked.west);
\node[label, text=TERed] at (7.05,-2.00) {Stop or replan};
\end{tikzpicture}
\caption{System overview.  Offline (left), a proposer supplies a candidate
effect descriptor, while counterfactual executions and policy-separating tests
determine whether it is refined, frozen, or sent for review.  Online (right),
the deterministic monitor instantiates the frozen descriptor for the exact
agent-proposed call, compares its effects with independently supplied authority,
and commits only an Allow{} call.}
\label{fig:system-overview}
\end{figure*}


Figure~\ref{fig:system-overview} separates two trust domains.  Offline
registration constructs and tests a tool descriptor against the registered
implementation.  Online mediation is deterministic: it instantiates the
frozen descriptor for a concrete call and supplies the resulting effect view
to a policy.  The LLM may seed candidate names and bindings during
registration, but neither its proposal nor a runtime agent decision is treated
as authority.

We use three terms consistently.  An \emph{effect atom} is the policy-relevant
unit in our model.  A \emph{tool descriptor} is the frozen output of offline
registration.  The \emph{provenance-origin monitor} is our AgentDojo runtime
instance: it expands descriptor fields into value records and rejects effects
introduced only by marked untrusted evidence.  A separate ToolSandbox
mechanism instantiates concrete typed atoms and checks them against an effect
bound.  ToolSandbox isolates the atom interface; AgentDojo evaluates a
conservative runtime instance over complete trajectories.

The registration evidence has two stages.  A committed-state difference is an
\emph{effect witness}: it shows that the intervention changes what the tool
does.  A pair becomes an \emph{authorization-separating witness} only when an
admissible policy decides the two resulting semantic views differently.  The
first justifies retaining information about an effect; the second justifies a
distinction in an authorization-sufficient representation.

\subsection{From Committed Effects to Policy Views}

Let $u=(t,\bar{x},s,h)$ contain a tool name, totalized arguments, pre-state, and
typed provenance evidence, and let $\mu(u)$ be the multiset of committed state changes produced by
executing $u$.  As in Section~\ref{sec:threat-model}, $\omega(u)$ augments this
multiset with policy-relevant trusted execution metadata.  An \emph{effect atom} records one occurrence that an admissible
policy may need to authorize independently.  For example, adding a participant
to an existing calendar event can produce
\[
 \langle\textsf{invite},\textsf{event:17},\textsf{sarah@example.com},
          \textsf{commit},\textsf{tool-output}\rangle.
\]
Here the components denote effect kind, operated resource, target principal,
operation or commit mode, and control provenance.  A tool-specific atom may
also retain a payload, visibility, date, or other qualifier when the policy
family can distinguish its values.  Lists expand per target or resource when
their members can be authorized independently.

A general registered descriptor $d_t$ contains (i) argument-totalization rules,
(ii) proposed effect kinds, (iii) effect-bearing schema fields and their
bindings or roles, (iv) inactive values, and (v) counterfactual evidence for
those choices.  Its instantiation function defines the monitor view
\[
  \rho_d(u)=A_{d_t}(u)
    =\operatorname{instantiate}(d_t,\bar{x},s,h),
\]
a canonical multiset of concrete atoms.  Thus $\mu(u)$ is what execution
actually commits, whereas $\rho_d(u)$ is the effect representation available
before commitment.  The ToolSandbox mechanism implements this typed form for
five source-executed tools.  The AgentDojo prototype implements the narrower
projection described above.  Its field roles are audit annotations rather than
independently validated semantic types: a field record is not automatically a
complete resource--principal--operation atom.

Atom equality is field-wise after descriptor-defined canonicalization, and
multiplicity is preserved: inviting the same principal to two resources is not
collapsed into one occurrence.  A general descriptor may bind aliases to stable
resource or principal identifiers.  The evaluated monitor has no such general
resolver.  Its frozen canonicalization handles bounded literal equivalences,
including case and whitespace, numeric and date forms, and normalized URL
forms.  A typed resolver may admit a value only from a matching authorized-read
ledger entry; it does not infer arbitrary entity aliases.  We therefore
evaluate alias-sensitive and typed-resource representations separately rather
than assuming canonicalization is semantically complete.

\paragraph{Policy-relative minimality.}
Minimality is relative to a frozen intervention domain $D$ and policy family
$\mathcal P$.  A representation is minimal on $(D,\mathcal P)$ when merging
any two of its classes creates a pair that some $P\in\mathcal P$ decides
differently.  Accordingly, changing committed state establishes that a field
is effect-relevant.  A trusted provenance or control-source change can also be
authorization-relevant without changing state.  Either establishes
authorization necessity only when an admissible policy separates the resulting
semantic views.  A returned-string change alone establishes neither.  This
definition does not assert a unique or globally minimal tuple schema.

\subsection{Offline Registration}

Registration consumes a versioned tool schema and implementation, sandbox
states, an intervention catalog, and, when authorization sufficiency is being
tested, a policy oracle from $\mathcal P$.  It proceeds as follows.

\paragraph{1. Propose.}
A developer or an LLM proposes effect names and field bindings from the tool
name, description, and schema.  The origin of the candidate is immaterial to
the security argument: it is an untrusted initial hypothesis and cannot
determine registration.  Our harness supports local-LLM seeding, but this paper
does not evaluate automatic contract synthesis or onboarding-time reduction.

\paragraph{2. Intervene and execute.}
For a valid base call, the harness changes a field value or its presence while
holding the implementation and initial state fixed.  Interventions include
typed alternatives, observed alternatives, boundary values, list expansion,
and omission/default cases when supported.  Surface-only variants serve as
placebos.  Both calls execute in copied sandboxes, so no external side effect
is produced.

\paragraph{3. Classify evidence.}
A state-difference oracle compares committed state, returned output, and
execution validity; provenance interventions additionally use the trusted
harness annotation.  The harness records \emph{committed-effect changed},
\emph{policy-metadata changed}, \emph{output-only changed},
\emph{semantic-invariant}, or \emph{invalid/unresolved}.  The first is evidence
that a field participates in the exercised effect relation; the second records
a change such as trusted control provenance.  Either becomes an
authorization-separating counterexample only if the selected policy oracle
distinguishes the two semantic views.  An invariant outcome supports omission
only over the tested states and interventions; an invalid outcome is never
treated as evidence of irrelevance.

\paragraph{4. Refine or freeze.}
A counterexample may add a field, split a list-valued effect, or introduce a
qualifier into the candidate representation.  Repeating this step implements
the finite refinement procedure analyzed in Section~\ref{sec:security-analysis}.
This is the general counterexample-guided loop; a descriptor is called
authorization-sufficient only after its test domain contains no
policy-separating collision.  The evaluated AgentDojo registry uses a more
conservative, nonminimal instance of this rule: it
retains fields with effect-change evidence and fields whose tests remain
unresolved.  A later multi-base audit broadens the evidence but does not
retroactively change that frozen registry.  Thus the registry demonstrates
coverage of exercised fields, not automatic discovery of a sparse minimum.
The descriptor, implementation identity, intervention manifest, and evidence
are hash-frozen together; runtime code cannot invoke the proposer.

\paragraph{Running example.}
Consider a calendar tool whose participant argument is a list.  The proposal
may initially describe the call only as \textsf{create-event}.  Registration
executes a base call that invites Alice and a paired call that additionally
invites Eve.  If the state oracle observes a second invitation and the policy
allows Alice but not Eve, the pair is an authorization-separating
counterexample.  Refinement therefore expands the participant list into one
\textsf{invite} atom per principal.  A surface paraphrase that preserves the
same event and participants should leave the atom multiset unchanged.  The
descriptor is frozen only for the intervention and policy domain exercised by
these checks; neither test proves correctness for arbitrary calendars or
future schema versions.

\subsection{Runtime Interface}

For context $u=(t,\bar{x},s,h)$, the runtime first applies the descriptor's
totalization function.  In the AgentDojo monitor, a versioned schema catalog
distinguishes required fields from explicit static defaults; inactive values
such as an omitted optional update remain represented but generate no active
effect check.  A runtime-derived value is accepted only when a registered
resolver relation and a matching typed authorized-read ledger entry justify
it.  Missing required values, unknown defaults, or unproved resolver values
produce \Abstain{} or a noncommitting replan decision.  The runtime then
instantiates $\rho_d(u)$ as defined above.  In the general architecture, a
policy interface
\[
  \operatorname{Authorize}(\rho_d(u),P)
       \rightarrow \{\Allow,\Deny,\Abstain\}
\]
receives these atoms and independently supplied authority context $P$.  Such a
policy may check ACLs, delegation, visibility, quotas, or other
resource-specific constraints.  Registration supplies the effect view; it
does not manufacture this authority.

The authority source is explicit in each evaluated setting.  Finite
representation validation enumerates bounds over observed source-effect
occurrences.  The concrete ToolSandbox mechanism treats one frozen
source-executed effect multiset as the bound and checks another call's typed
atoms for multiset inclusion.  AgentDojo supplies no general ACL: its bound is
the narrower provenance-origin relation between the authenticated task,
registered effects, and typed untrusted evidence.  These settings test the
representation and mediation interfaces under stated policies; they are not
interchangeable evidence for deployment authority.

\subsection{Provenance-Origin Monitor}

The provenance-origin monitor instantiates a narrower policy than the general
interface above.  It asks whether explicitly untrusted observations introduced
a registered effect or effect-bearing value absent from the authenticated task.

\paragraph{Structured provenance.}
The adapter traverses mappings, sequences, dataclasses, and Pydantic-like
records before extracting explicitly marked untrusted control segments.  This
preserves provenance through structured tool outputs.  Provenance identifies
where a value came from; it does not by itself authorize the resulting effect.

\paragraph{Registered-field binding.}
For each registered field, the monitor expands mappings and lists to concrete
leaf values.  A value is rejected when it occurs in untrusted control text but
is not literally grounded in the authenticated task.  Per-value expansion
prevents one task-grounded recipient from authorizing an additional injected
recipient.  The audit record stores the field, concrete value, descriptor role,
provenance status, and decision reason.

\paragraph{Effect-label grounding.}
Some injected requests change the intended action without exposing a unique
argument value.  The monitor therefore performs a second, lexical check over action and
object terms derived from the frozen tool name and effect label.  If the
untrusted segment requests that registered effect while the authenticated task
does not, the call is denied.  This is a benchmark-specific grounding heuristic,
not general semantic authorization or concrete atom recovery.

Missing registration or failed totalization produces \Abstain; a detected
untrusted expansion produces \Deny; otherwise the monitor returns \Allow.  In
the calendar example, an Eve value copied from explicitly untrusted text is
denied when Eve is absent from the authenticated task; the monitor does not
claim to infer whether the user has organization-level permission to invite her.
The executor
commits only the exact totalized call signature associated with \Allow, and the
audit log reconciles checked and executed calls.  The monitor makes no runtime
LLM call.  ACLs and other authority predicates remain external to this
evaluated policy instance.
